\documentclass[10pt,a4paper]{amsart}
\usepackage{amsmath,amssymb,amsthm,graphicx}
% Retained minimal preamble for the Stacks Project chapter more-algebra.tex (More on Algebra).
% The chapter file pinned-original-source.tex carries only \input{preamble} on its first line, so no
% preamble text lives inside the chapter itself.  The Stacks Project uses one single shared preamble for
% every chapter; the snapshot stages it identically as sources/stacks-chow/preamble.tex and
% sources/stacks-etale-cohomology/preamble.tex.  The block below is copied mechanically, line for line,
% from sources/stacks-chow/preamble.tex lines 14-18 (the Xy-pic setup) and 154-217 (the native theorem
% declarations together with the mathematical macros).  Nothing was added, removed or redefined; the
% wrapper class and the amsmath/amssymb/amsthm/graphicx base packages are supplied by the assembler.
% Chapter: more-algebra.tex -- More on Algebra.
% For commutative diagrams we use Xy-pic
\usepackage[all]{xy}

% We use 2cell for 2-commutative diagrams.
\xyoption{2cell}

% Theorem environments.
%
\theoremstyle{plain}
\newtheorem{theorem}[subsection]{Theorem}
\newtheorem{proposition}[subsection]{Proposition}
\newtheorem{lemma}[subsection]{Lemma}

\theoremstyle{definition}
\newtheorem{definition}[subsection]{Definition}
\newtheorem{example}[subsection]{Example}
\newtheorem{exercise}[subsection]{Exercise}
\newtheorem{situation}[subsection]{Situation}

\theoremstyle{remark}
\newtheorem{remark}[subsection]{Remark}
\newtheorem{remarks}[subsection]{Remarks}

\numberwithin{equation}{subsection}

% Macros
%
\def\lim{\mathop{\mathrm{lim}}\nolimits}
\def\colim{\mathop{\mathrm{colim}}\nolimits}
\def\Spec{\mathop{\mathrm{Spec}}}
\def\Hom{\mathop{\mathrm{Hom}}\nolimits}
\def\Ext{\mathop{\mathrm{Ext}}\nolimits}
\def\SheafHom{\mathop{\mathcal{H}\!\mathit{om}}\nolimits}
\def\SheafExt{\mathop{\mathcal{E}\!\mathit{xt}}\nolimits}
\def\Sch{\mathit{Sch}}
\def\Mor{\mathop{\mathrm{Mor}}\nolimits}
\def\Ob{\mathop{\mathrm{Ob}}\nolimits}
\def\Sh{\mathop{\mathit{Sh}}\nolimits}
\def\NL{\mathop{N\!L}\nolimits}
\def\CH{\mathop{\mathrm{CH}}\nolimits}
\def\proetale{{pro\text{-}\acute{e}tale}}
\def\etale{{\acute{e}tale}}
\def\QCoh{\mathit{QCoh}}
\def\Ker{\mathop{\mathrm{Ker}}}
\def\Im{\mathop{\mathrm{Im}}}
\def\Coker{\mathop{\mathrm{Coker}}}
\def\Coim{\mathop{\mathrm{Coim}}}

% Boxtimes
%
\DeclareMathSymbol{\boxtimes}{\mathbin}{AMSa}{"02}

%
% Macros for moduli stacks/spaces
%
\def\QCohstack{\mathcal{QC}\!\mathit{oh}}
\def\Cohstack{\mathcal{C}\!\mathit{oh}}
\def\Spacesstack{\mathcal{S}\!\mathit{paces}}
\def\Quotfunctor{\mathrm{Quot}}
\def\Hilbfunctor{\mathrm{Hilb}}
\def\Curvesstack{\mathcal{C}\!\mathit{urves}}
\def\Polarizedstack{\mathcal{P}\!\mathit{olarized}}
\def\Complexesstack{\mathcal{C}\!\mathit{omplexes}}
% \Pic is the operator that assigns to X its picard group, usage \Pic(X)
% \Picardstack_{X/B} denotes the Picard stack of X over B
% \Picardfunctor_{X/B} denotes the Picard functor of X over B
\def\Pic{\mathop{\mathrm{Pic}}\nolimits}
\def\Picardstack{\mathcal{P}\!\mathit{ic}}
\def\Picardfunctor{\mathrm{Pic}}
\def\Deformationcategory{\mathcal{D}\!\mathit{ef}}

% --- appended 2026-10-12 (Chief): the project own definitions of the reference/slogan/history
% display environments, copied mechanically from the shared Stacks preamble lines 8-12
% (\usepackage{verbatim} plus the three \newenvironment lines, which the chapter preamble keeps
% verbatim from the project preamble). Nothing added, removed or redefined: these environments
% carry no mathematical content and the project itself typesets them as comments. Without these
% three lines a counted statement whose own text uses them fails with Environment slogan undefined.
\usepackage{verbatim}
\newenvironment{reference}{\comment}{\endcomment}
\newenvironment{slogan}{\comment}{\endcomment}
\newenvironment{history}{\comment}{\endcomment}
\usepackage{cleveref}
\crefname{originallemma}{original lemma}{original lemmas}
\Crefname{originallemma}{Original lemma}{Original lemmas}
\makeatletter
\newlabel{more-algebra-lemma-sequence-reflexive}{{0EB8}{source}{}{sourcepointer-1}{}}
\newlabel{more-algebra-lemma-sequence-reflexive@cref}{{[originallemma][0][]0EB8}{[1][0][]source}{}{}{}}
\makeatother
\crefname{originallemma}{original lemma}{original lemmas}
\Crefname{originallemma}{Original lemma}{Original lemmas}
\makeatletter
\newlabel{algebra-lemma-ext-noetherian}{{08YR}{source}{}{sourcepointer-2}{}}
\newlabel{algebra-lemma-ext-noetherian@cref}{{[originallemma][0][]08YR}{[1][0][]source}{}{}{}}
\makeatother
\crefname{originallemma}{original lemma}{original lemmas}
\Crefname{originallemma}{Original lemma}{Original lemmas}
\makeatletter
\newlabel{lemma-reflexive-torsion-free}{{0AV0}{source}{}{sourcepointer-3}{}}
\newlabel{lemma-reflexive-torsion-free@cref}{{[originallemma][0][]0AV0}{[1][0][]source}{}{}{}}
\makeatother
\crefname{originallemma}{original lemma}{original lemmas}
\Crefname{originallemma}{Original lemma}{Original lemmas}
\makeatletter
\newlabel{algebra-lemma-hom-from-finitely-presented}{{0583}{source}{}{sourcepointer-4}{}}
\newlabel{algebra-lemma-hom-from-finitely-presented@cref}{{[originallemma][0][]0583}{[1][0][]source}{}{}{}}
\makeatother
\title{Reflexivity of the kernel of a surjection with reflexive middle term and torsion-free quotient}
\author{The Stacks Project Authors}
\date{stacks-project:more-algebra.tex:local-snapshot-20261012; GNU FDL 1.2 or later}
\begin{document}
\maketitle
\setcounter{subsection}{1}
\noindent\textit{Editorial note.} All delivered mathematical text is the author's own, copied from the pinned Stacks Project chapter file more-algebra.tex; no statement or proof step was rewritten, reordered or removed.\par\medskip
\noindent\textit{Editorial note.} The chapter file pinned-original-source.tex carries only \textbackslash{}input\{preamble\} on its first line, so no preamble text lives inside the chapter. The retained file sources/stacks-more-algebra/retained-preamble.tex is a mechanical, line-for-line copy of the project's single shared preamble as staged in the snapshot (Xy-pic setup plus theorem declarations and mathematical macros). Nothing was added, removed or redefined, and the excerpt serves only as preamble, never as a span coordinate.\par\medskip
\noindent\textit{Editorial note.} Both counted spans start at the \textbackslash{}begin line of their environment, so the statement and the author's proof are delivered whole; surrounding prose of the section is deliberately outside the counted spans and is carried as standard prerequisites instead.\par\medskip
\noindent\textit{Editorial note.} The citation fields (source\_id, source\_version, authors, source\_url, license.name, source\_title) are identical to those of every other card this lane has produced from this chapter, so the batch selector treats them as one source.\par\medskip
\noindent\textit{Allowed standard prerequisites (editorial).}\begin{itemize}
\item Hom modules between finite modules over a Noetherian ring are finite. Algebra, Lemma \textbackslash{}ref\{algebra-lemma-ext-noetherian\}, cited at proof lines 5263-5264. Makes the duals and double duals finite modules, so the diagram chase is a statement about finite modules.
\item A finite module over a Noetherian domain is torsion free if and only if the map to its double dual is injective. Lemma \textbackslash{}ref\{lemma-reflexive-torsion-free\}, cited at proof line 5272. Gives injectivity of the right vertical arrow of the diagram.
\item Localisation at the fraction field is compatible with Hom out of a finitely presented module. Algebra, Lemma \textbackslash{}ref\{algebra-lemma-hom-from-finitely-presented\}, cited at proof line 5288. Produces the exact sequence of vector spaces whose exactness forces Q tensored with the fraction field to vanish.
\end{itemize}\medskip
\ifdefined\hypertarget\hypertarget{sourcepointer-1}{}\fi
\noindent\textit{Source reference.} Original lemma 0EB8: focused-seeds/AGP/sources/stacks-snapshot-20261011/tags.txt, source lines 17913--17913. The lemma's own canonical tag, quoted so that readers can resolve the label inside the delivered text; it points at the counted statement itself and carries no proof content.\par\medskip
\ifdefined\hypertarget\hypertarget{sourcepointer-2}{}\fi
\noindent\textit{Source reference.} Original lemma 08YR: focused-seeds/AGP/sources/stacks-snapshot-20261011/tags.txt, source lines 11349--11349. Navigation pointer to a statement the counted proof refers to but does not contain; nothing load-bearing is delegated to it.\par\medskip
\ifdefined\hypertarget\hypertarget{sourcepointer-3}{}\fi
\noindent\textit{Source reference.} Original lemma 0AV0: focused-seeds/AGP/sources/stacks-snapshot-20261011/tags.txt, source lines 13668--13668. Navigation pointer to a statement the counted proof refers to but does not contain; nothing load-bearing is delegated to it.\par\medskip
\ifdefined\hypertarget\hypertarget{sourcepointer-4}{}\fi
\noindent\textit{Source reference.} Original lemma 0583: focused-seeds/AGP/sources/stacks-snapshot-20261011/tags.txt, source lines 6766--6766. Navigation pointer to a statement the counted proof refers to but does not contain; nothing load-bearing is delegated to it.\par\medskip
\begin{lemma}
\label{lemma-sequence-reflexive}
Let $R$ be a Noetherian domain. Let $0 \to M \to M' \to M''$
an exact sequence of finite $R$-modules. If $M'$ is reflexive
and $M''$ is torsion free, then $M$ is reflexive.
\end{lemma}

\begin{proof}
We will use without further mention that $\Hom_R(N, N')$ is a finite
$R$-module for any finite $R$-modules $N$ and $N'$, see
Algebra, Lemma \ref{algebra-lemma-ext-noetherian}.
We take duals to get a sequence
$$
\Hom_R(M, R) \leftarrow \Hom_R(M', R) \leftarrow \Hom_R(M'', R)
$$
Dualizing again we obtain a commutative diagram
$$
\xymatrix{
\Hom_R(\Hom_R(M, R), R) \ar[r]_j &
\Hom_R(\Hom_R(M', R), R) \ar[r] &
\Hom_R(\Hom_R(M'', R), R) \\
M \ar[u] \ar[r] & M' \ar[u] \ar[r] & M'' \ar[u]
}
$$
We do not know the top row is exact. But, by assumption
the middle vertical arrow is an isomorphism and the right
vertical arrow is injective (Lemma \ref{lemma-reflexive-torsion-free}).
We claim $j$ is injective. Assuming the claim
a diagram chase shows that the left vertical
arrow is an isomorphism, i.e., $M$ is reflexive.

\medskip\noindent
Proof of the claim. Consider the exact sequence
$\Hom_R(M',R)\to \Hom_R(M,R)\to Q \to 0$
defining $Q$. One applies
Algebra, Lemma \ref{algebra-lemma-hom-from-finitely-presented}
to obtain
$$
\Hom_K(M'\otimes_R K,K) \to
\Hom_K(M\otimes_R K,K) \to
Q\otimes_R K\to 0
$$
But $M \otimes_R K \to M' \otimes_R K$ is an injective map of vector spaces,
hence split injective, so $Q \otimes_R K = 0$, that is, $Q$ is torsion.
Then one gets the exact sequence 
$$
0 \to \Hom_R(Q,R) \to
\Hom_R(\Hom_R(M,R),R) \to
\Hom_R(\Hom_R(M',R),R)
$$
and $\Hom_R(Q,R)=0$ because $Q$ is torsion.
\end{proof}

\end{document}
